% ===== 问题一正文 =====
% 在本子文件修改问题一正文，公式和图表编号由主文件统一管理。
\subsection{问题一：交会定位法计算多边形定位区域直径}\label{sec:q1}
\subsubsection{计算干扰源定位区域}\label{sec:q1-1}

设机器狗在 $n$ 个位置 $S_i=(x_i,y_i)$ 处对该干扰源测得示向度 $\theta_i$。$\theta_i$ 是自正东逆时针量到“检测点指向干扰源”方向的角度，但它并不是干扰源的真实方位：真实方位落在 $\theta_i$ 两侧各 $1^\circ$ 内。因此，干扰源一定位于以 $S_i$ 为顶点、方向在 $\theta_i\pm1^\circ$ 之间的一个 $2^\circ$ 窄楔形内。

现在要把“点在楔形内”写成数学不等式。记 $\varphi_i=\pi\theta_i/180$，$\delta=\pi\varepsilon/180=\pi/180$，楔形的两条边方向分别为

\[
u_i^-=\begin{pmatrix}\cos(\varphi_i-\delta)\\ \sin(\varphi_i-\delta)\end{pmatrix},\qquad
u_i^+=\begin{pmatrix}\cos(\varphi_i+\delta)\\ \sin(\varphi_i+\delta)\end{pmatrix},
\]

其中 $u_i^-$ 对应偏转角度的下限 $\varphi_i-\delta$，$u_i^+$ 对应上限 $\varphi_i+\delta$。判断一个点在两条边的哪一侧，可以用二维向量的叉积：对 $a=(a_x,a_y)$、$b=(b_x,b_y)$，定义 $a\times b=a_xb_y-a_yb_x$，它大于零说明 $b$ 在 $a$ 的逆时针一侧，小于零说明在顺时针一侧。于是，点 $X=(x,y)$ 落在楔形内的条件是

\begin{equation}\label{eq:q1-1}
u_i^-\times(X-S_i)\ge 0,\qquad u_i^+\times(X-S_i)\le 0.
\end{equation}

第一式要求 $X-S_i$ 不越过下边 $u_i^-$（在其逆时针一侧），第二式要求不越过上边 $u_i^+$（在其顺时针一侧），两式合起来就是把 $X$ 夹在两条边之间。

将式\eqref{eq:q1-1} 展开：

\begin{equation}\label{eq:q1-2}
\begin{aligned}
\sin(\varphi_i-\delta)(x-x_i)-\cos(\varphi_i-\delta)(y-y_i)&\le 0,\\
-\sin(\varphi_i+\delta)(x-x_i)+\cos(\varphi_i+\delta)(y-y_i)&\le 0.
\end{aligned}
\end{equation}

这两个式子都是 $x$、$y$ 的一次不等式，形如 $ax+by\le c$。满足 $ax+by\le c$ 的点集是直线 $ax+by=c$ 一侧的半个平面，称为半平面，它是一个凸集。式\eqref{eq:q1-2} 的两个不等式分别对应楔形的两条边，它们的交集就是该楔形。这样就把几何上“点在楔形内”的条件，转化成了代数上两个线性不等式同时成立。

$n$ 次测向后，干扰源必须同时落在每个楔形内，也就是要同时满足全部 $2n$ 个不等式。将它们统一写成 $a_jx+b_jy\le c_j$，得到定位区域

\begin{equation}\label{eq:q1-3}
P=\bigcap_{j=1}^{2n}\left\{(x,y):a_jx+b_jy\le c_j\right\}.
\end{equation}

$P$ 是 $2n$ 个半平面的交。半平面是凸集，凸集之交仍是凸集，故 $P$ 为凸集；在非空且有界时，$P$ 是\textbf{凸多边形}，退化情形为线段或点。题目图 2 中的红色四边形就是 $n=2$ 时两个楔形求交的结果。每增加一次有效测向，约束增加两条，定位区域随之收缩。

\subsubsection{求定位区域的直径}\label{sec:q1-2}

定位区域 $P$ 是凸多边形，现在要求它的直径---区域内距离最远的一对点之间的距离。

\textbf{首先确认区域有界。} 并非任意一组测向都能定出有限的区域。若各检测点近似共线、示向度方向接近平行，楔形之交沿某方向无限延伸，区域无界，直径为无穷大；若某次读数偏差较大，各楔形可能互不相容，区域为空。这两种情形都需要先排除。

具体做法是用线性规划：先检验约束组 $\{a_jx+b_jy\le c_j\}$ 是否可行，不可行说明数据矛盾，应剔除异常的测向数据；若可行，再分别求 $x$、$-x$、$y$、$-y$ 在约束下的最大值，四者均有限当且仅当 $P$ 有界。注意，无界区域也可能有若干有限顶点（例如楔形的尖端），若跳过这一步直接取顶点对距离，会把本应为无穷的直径误报成有限值。

\textbf{确认有界后，求出全部顶点。} 凸多边形的每个顶点都是某两条边界直线的交点。对直线 $a_jx+b_jy=c_j$ 与 $a_kx+b_ky=c_k$，记 $\Delta_{jk}=a_jb_k-a_kb_j$，当 $\Delta_{jk}\ne0$ 时交点为

\begin{equation}\label{eq:q1-4}
x_{jk}=\frac{c_jb_k-b_jc_k}{\Delta_{jk}},\qquad
y_{jk}=\frac{a_jc_k-c_ja_k}{\Delta_{jk}}.
\end{equation}

$\Delta_{jk}=0$ 时两直线平行，没有孤立交点，跳过。求出的交点未必在区域内，需代回全部 $2n$ 个约束检验，保留满足所有不等式者；去重后即得顶点集 $\{v_1,\dots,v_h\}$。

\textbf{有了顶点，直径就可以确定。} 按题目定义，$D=\max_{X,Y\in P}\|X-Y\|_2$。因为 $P$ 是凸多边形，我们证明这个最大值一定在顶点对上取得：

\begin{equation}\label{eq:q1-5}
D=\max_{1\le i<j\le h}\|v_i-v_j\|_2 .
\end{equation}

理由如下。凸多边形是其全部顶点的凸包，因此 $P$ 中任意一点都可写为顶点的凸组合：存在 $\lambda_i\ge0$、$\sum_i\lambda_i=1$，使 $X=\sum_i\lambda_iv_i$；同理 $Y=\sum_j\mu_jv_j$，$\mu_j\ge0$、$\sum_j\mu_j=1$。两式相减得

\[
X-Y=\sum_{i,j}\lambda_i\mu_j(v_i-v_j),
\qquad
\sum_{i,j}\lambda_i\mu_j=\left(\sum_i\lambda_i\right)\left(\sum_j\mu_j\right)=1 .
\]

由三角不等式，

\[
\|X-Y\|\le\sum_{i,j}\lambda_i\mu_j\|v_i-v_j\|\le\max_{i,j}\|v_i-v_j\|\cdot1=\max_{i,j}\|v_i-v_j\|,
\]

即区域内任意两点的距离都不超过最远顶点对的距离。而最远顶点对本身就是区域内的点，所以这个上界可以取到。式\eqref{eq:q1-5} 成立。

综上，求直径的完整流程为：

\par\medskip
\noindent\begin{minipage}{\linewidth}
\small
\refstepcounter{algorithm}\label{alg:q1-diameter}
\textbf{算法\thealgorithm：定位区域直径计算}\par\smallskip
\noindent\textbf{输入：}检测点 $S_i$、示向度 $\theta_i$，$i=1,\dots,n$
\begin{enumerate}
\setlength{\itemsep}{2pt}
\setlength{\parskip}{0pt}
\item 按式\eqref{eq:q1-2} 将每次测向写成两个半平面约束；
\item 线性规划检验可行性与有界性；
\item 按式\eqref{eq:q1-4} 求每对边界直线的交点；
\item 保留满足全部约束的交点并去重，得顶点集；
\item 比较所有顶点对的距离，取最大者 $(A,B)$，$D=|AB|$。
\end{enumerate}
\end{minipage}
\par\medskip

边界直线共 $2n$ 条，交点约 $\binom{2n}{2}=O(n^2)$ 个，每个交点检验约束需 $O(n)$，顶点数不超过 $2n$，比较顶点对为 $O(n^2)$，总计算量 $O(n^3)$。本题 $n$ 一般不超过十几，计算量完全可以实时完成。

\subsubsection{判断直径圆能否覆盖定位区域}\label{sec:q1-3}

机器狗的光学探测仪只能在距干扰源 $20$ m 以内工作。如果以定位区域直径 $D$ 为直径画一个圆，能盖住整个区域，那么只要 $D\le40$ m，机器狗走到圆心就能保证清除成功---这比计算最小包围圆要简单得多。所以问题是：这样的圆真的能盖住整个区域吗？

首先要确定圆心。设最远的一对顶点为 $A$、$B$，距离为 $D$。若半径 $D/2$、圆心 $C$ 的圆覆盖了 $P$，则 $A$、$B$ 都在圆内，即 $\|A-C\|\le D/2$、$\|C-B\|\le D/2$。代入三角不等式

\[
D=\|A-B\|\le\|A-C\|+\|C-B\|\le\frac D2+\frac D2=D .
\]

不等式的两端相等，故中间每一步都取等号：$\|A-C\|=\|C-B\|=D/2$，且 $C$ 在线段 $AB$ 上。这两个条件合起来，$C$ 只能是 $AB$ 的中点 $C_0=(A+B)/2$。圆是凸集，覆盖全部顶点就覆盖了整个凸多边形，所以判据为

\begin{equation}\label{eq:q1-6}
\max_{1\le i\le h}\|v_i-C_0\|\le\frac D2 .
\end{equation}

式\eqref{eq:q1-6} 是否一定满足？下面说明它可以不成立，并给出一个满足题设条件的具体反例。

\textbf{构造思路：} 要让直径圆盖不住区域，需要找到这样一个顶点：它已经跑到了直径圆外面，但它到其他顶点的距离还没有超过原来的直径---直径没变，圆却已经不够大了。

先用归一化的坐标把这个想法量化。取 $A=(-1,0)$、$B=(1,0)$，$|AB|=2$，直径圆以原点为圆心、半径为 $1$。在上方放一个点 $U=(0,h)$，希望它同时满足：

\begin{itemize}
\item 跑到圆外：$h>1$；
\item 不改变直径：$|UA|=|UB|=\sqrt{1+h^2}\le2$，即 $h\le\sqrt{3}$。
\end{itemize}

两个条件合起来是 $1<h\le\sqrt{3}$，这个区间非空。直观地理解：上顶点刚越过圆的最高点时，它离左右两端还不够远，不会产生更长的边。

取 $h=1.5$，则 $|UA|=\sqrt{3.25}\approx1.803<2$，直径仍为 $2$；而上顶点距圆心 $1.5$，超出半径 $50\%$。反例的形状就有了：底边为 $2$、高为 $1.5$ 的等腰三角形。

\textbf{用测向数据实现该形状：} 将坐标放大到米，目标三角形的三个顶点为

\[
A=(-10,\,0),\quad B=(10,\,0),\quad U=(0,\,15).
\]

三条边分别是 $y=0$（底边）、$y=1.5x+15$（左腰）、$y=-1.5x+15$（右腰）。三角形有 $3$ 条边，而每个检测点的楔形提供 $2$ 条边界线，因此用 $3$ 个检测点（$6$ 条边界线）即可实现：让其中 $3$ 条边界线分别与三角形的三条边重合，另外 $3$ 条不切入三角形内部。具体取法见表\ref{tab:q1-counterexample}。

\begin{center}
\begin{minipage}{\linewidth}
\centering
\captionof{table}{反例的检测数据（误差界 $\pm1^\circ$）}\label{tab:q1-counterexample}

\renewcommand{\arraystretch}{1.3}
\begin{tabularx}{0.98\linewidth}{>{\centering\arraybackslash}p{0.22\linewidth}>{\centering\arraybackslash}p{0.20\linewidth}>{\raggedright\arraybackslash}X}
\toprule
\textbf{检测点 $S_i$（m）} & \textbf{示向度 $\theta_i$} & \textbf{起决定作用的边界} \\
\midrule
$(-500,\,0)$ & $1^\circ$ & 下边界方向为 $0^\circ$，给出 $y\ge0$（底边） \\
$(-310,\,-450)$ & $\approx55.31^\circ$ & 上边界方向为 $\arctan(1.5)$，给出 $y\le1.5x+15$（左腰） \\
$(310,\,-450)$ & $\approx124.69^\circ$ & 下边界方向为 $180^\circ-\arctan(1.5)$，给出 $y\le-1.5x+15$（右腰） \\
\bottomrule
\end{tabularx}
\end{minipage}
\end{center}

其中第二个检测点的示向度取 $\theta_2=\arctan(1.5)-1^\circ\approx55.31^\circ$，第三个取 $\theta_3=181^\circ-\arctan(1.5)\approx124.69^\circ$。将三角形的三个顶点分别代入每个楔形的另一条边界线，可以验证它们均满足约束（不被裁掉），因此 $6$ 个半平面的交恰好就是三角形 $ABU$。

现在检验覆盖性。边长为 $|AB|=20$，$|AU|=|BU|=\sqrt{325}\approx18.03$，直径 $D=20$ m 由 $A$、$B$ 取得。直径圆圆心 $C_0=(0,\,0)$，半径 $D/2=10$ m。而上顶点到圆心的距离为

\[
\|U-C_0\|=15\ \text{m}>10\ \text{m}=\frac D2,
\]

超出半径 $5$ m，式\eqref{eq:q1-6} 不成立。

\begin{center}
\begin{minipage}{\linewidth}
\centering
\QOneCaseFigure
\captionof{figure}{三次测向的角域交集与直径圆覆盖检验。左图按实际比例绘制三个误差角域；右图放大定位三角形，红色圆的半径为 $10$ m，上顶点超出 $5$ m。}
\label{fig:q1-counterexample}
\end{minipage}
\end{center}

三个检测点到三角形中心的距离均约 $500$ m，在有效接收半径内；三角形整体位于半径 $1800$ m 的目标区域内，是完全合理的探测场景。于是问题一的第二问答案为否：

\textbf{以定位区域直径为直径的圆，不一定能覆盖该定位区域。}


